\subsection{Background and motivation}

The deployment of IoT sensor networks has garnered significant research attention due to
its wide-ranging applications in smart agriculture, environmental monitoring, and disaster
response~\cite{wei2022uav}. In many of these scenarios, fixed infrastructure is either
unavailable or prohibitively expensive to deploy. Consequently, UAVs have emerged as a
critical enabler for data collection, offering the flexibility to serve as mobile gateways
that can be deployed on-demand to remote or hazardous locations~\cite{wei2022uav,mozaffari2019tutorial,zeng2016wireless}. By
combining the long-range, low-power capabilities of the LoRaWAN protocol~\cite{augustin2016lora,bor2016lora}
with the mobility of UAVs, it is theoretically possible to collect data from vast sensor
fields with minimal infrastructure~\cite{nakamura2022lora}.

However, integrating mobile UAVs into dense LoRaWAN networks presents fundamental challenges
that standard protocols were not designed to handle. The LoRaWAN architecture, specifically
its ADR mechanism, is engineered for static, ground-based networks. As identified by Zorbas
et al.\ (2025), in dense deployments where link qualities change rapidly---such as with a
moving UAV---the standard network management fails. Devices often converge to the same
transmission settings, leading to a high rate of ``intra-SF collisions'' and a degradation
of network performance~\cite{zorbas2025revisiting}. This phenomenon, often referred to as
the ``SF12 Well'' or ``SF Monoculture,'' causes a congestion collapse where the network's
concurrent channel capacity is wasted, and packet delivery rates plummet.

Existing solutions to this problem remain siloed. Standard path planning heuristics (such as
``lawnmower'' patterns) optimise solely for flight distance, ignoring the underlying
telecommunications physics~\cite{xiong2023flyinglora}. These ``protocol-blind'' trajectories
inadvertently force sensors into the SF Monoculture, causing the UAV to hover inefficiently
while waiting for data that cannot be decoded due to collisions. Conversely, mathematical
optimisation approaches (like Mixed Integer Linear Programming) can solve the resource
allocation problem for static networks but lack the real-time adaptability required for
dynamic, mobile environments~\cite{zorbas2025revisiting}.

To address these limitations, a protocol-aware control framework is required---one that
integrates the UAV's physical mobility with the logical constraints of the LoRaWAN protocol.
Recent literature has explicitly suggested that reinforcement learning models could be employed
to adapt to these evolving network conditions~\cite{zorbas2025revisiting,nguyen20226g,luong2019drl,challita2018interference}. By
treating the UAV's 3D position as a direct control input for the radio protocol, an intelligent
agent can proactively manage the RF environment, breaking the SF monoculture and ensuring
true energy efficiency through optimised network throughput.
